% PAGES: 286-286
% sec09_c2_7.tex ends with a fully-closed \end{align} for (9.66) (the comma
% printed at the end of that align block is grammatical, not a LaTeX
% construct left open): the compound sentence "we find that [[align]] ...
% w_T, since, from (9.60), it follows that [[display]] ... = 1-w_{max,cash}."
% splits across the page/file boundary between its \end{align} and the fresh
% \[ display below. No environment is left open across that boundary. A new
% \begin{answerx} opens near the bottom of this file and IS left OPEN at the
% end -- it continues on PDF page 287 (folio 271), which closes it with the
% printed square. sec09_c2_9.tex must NOT reopen it.
%
% SOURCE INCONSISTENCY: "the formulas (9.65) and (9.66)" below is printed
% exactly that way, even though the two bulleted equations immediately above
% it in this same paragraph are numbered (9.67) and (9.68), not (9.65)/(9.66)
% (which are the unbulleted, page-285 versions of the same two formulas). The
% analogous sentence in section 9.3 ("...computed using the formulas (9.40)
% and (9.41)...", sec09_c2_1.tex) correctly cites the directly-preceding
% bulleted pair, so this looks like a genuine numbering slip in the source,
% not a deliberate cross-reference. Transcribed as printed.

since, from (9.60), it follows that
\[
\sigma_P\frac{\bfone^t\Sigma_R^{-1}\bar\mu}{\sqrt{\bar\mu^t\Sigma_R^{-1}\bar\mu}} \;=\;
1-w_{max,cash}.
\]

From (9.65) and (9.66), we conclude that the maximum expected return of a portfolio
with variance of return equal to $\sigma_P^2$ is obtained for the following asset
allocation:
\begin{itemize}
\item weight $w_{max,cash}$ of the cash position given by
\begin{equation}
w_{max,cash} \;=\; 1-\frac{\sigma_P}{\sqrt{w_T^t\Sigma_Rw_T}}\cdot
\operatorname{sign}\!\left(\bfone^t\Sigma_R^{-1}\bar\mu\right);
\end{equation}
\item asset weights vector $w_{max}$ given by
\begin{equation}
w_{max} \;=\; (1-w_{max,cash})\,w_T.
\end{equation}
\end{itemize}

The pseudocode for the asset allocation for the maximum return portfolio computed
using the formulas (9.65) and (9.66) can be found in Table 9.4.

We revisit below the example from section 9.4 and find the maximum return portfolio
by using the tangency portfolio.

\par\medskip\noindent\textit{Example:}\ Find the \$100 million maximum return
portfolio with 27\% standard deviation of return invested in cash and four assets
with expected returns equal to 2\%, 1.75\%, 2.5\%, and 1.5\%, and with the following
covariance matrix of their returns:
\[
\begin{pmatrix}
0.09 & 0.01 & 0.03 & -0.015 \\
0.01 & 0.0625 & -0.02 & -0.01 \\
0.03 & -0.02 & 0.1225 & 0.02 \\
-0.015 & -0.01 & 0.02 & 0.0576
\end{pmatrix}.
\]

Assume that the risk free rate is 1\%.

\begin{answerx}
% Left OPEN at the end of this file -- closes on PDF page 287 (folio 271).
Let $\mu_1=0.02$, $\mu_2=0.0175$, $\mu_3=0.025$, $\mu_4=0.015$, and $r_f=0.01$. We
look for the asset and cash allocation corresponding to the maximum return portfolio
with standard deviation of return equal to $\sigma_P = 0.27$.

Note that $\bar\mu=\mu-r_f\bfone$ and $\Sigma_R$ are given by
\[
\bar\mu \;=\; \begin{pmatrix} 0.01 \\ 0.0075 \\ 0.015 \\ 0.005 \end{pmatrix}; \quad
\Sigma_R \;=\; \begin{pmatrix}
0.09 & 0.01 & 0.03 & -0.015 \\
0.01 & 0.0625 & -0.02 & -0.01 \\
0.03 & -0.02 & 0.1225 & 0.02 \\
-0.015 & -0.01 & 0.02 & 0.0576
\end{pmatrix}.
\]

The vector $\Sigma_R^{-1}\bar\mu = \text{linear\_solve\_cholesky}(\Sigma_R,\bar\mu)$ is
computed by using the Cholesky decomposition to solve the linear system corresponding
to the matrix $\Sigma_R$ and to the right hand side $\bar\mu$.

From (9.62), we find that the asset weights vector of the tangency portfolio is
\[
w_T \;=\; \frac{1}{\bfone^t\Sigma_R^{-1}\bar\mu}\Sigma_R^{-1}\bar\mu \;=\;
\begin{pmatrix} 0.1587 \\ 0.3660 \\ 0.2650 \\ 0.2103 \end{pmatrix}.
\]
